% Source: Weinberg GR, physical PDF pp. 146--149
%         (printed pp. 121--124).
% Coverage: Section 5.1, Particle Mechanics; equations (5.1.1)--(5.1.14).
% Figures/tables/footnotes: no figures or tables; reference markers 1 and 2.
% Status: source-reviewed and compile-clean.
% Uncertainties: none.

\subsection{Particle Mechanics}\label{sec:5.1}

A particle not under the influence of any force will in special relativity
have constant four-velocity \(u^a\) and constant spin \(S_a\), that is,
\begin{equation}
  \frac{\dd u^a}{\dd\tau}=0,
  \qquad
  \left(
    u^a\equiv\frac{\dd\xi^a}{\dd\tau}
  \right),
  \tag{5.1.1}\label{eq:5.1.1}
\end{equation}
\begin{equation}
  \frac{\dd S_a}{\dd\tau}=0.
  \tag{5.1.2}\label{eq:5.1.2}
\end{equation}
Recall that \(S_a\) is defined in the rest frame of the particle to have the
time-first components \(\{0,\mathbf S\}\), so that in an arbitrary Lorentz
frame it satisfies the further relation
\begin{equation}
  S_a u^a=0.
  \tag{5.1.3}\label{eq:5.1.3}
\end{equation}

In order to make these equations generally covariant, we define vectors
\(u^\mu\) and \(S_\mu\) in a general coordinate system \(x^\mu\) by
\begin{equation}
  u^\mu
  \equiv
  \frac{\partial x^\mu}{\partial\xi^a}u^a
  =
  \frac{\dd x^\mu}{\dd\tau},
  \tag{5.1.4}\label{eq:5.1.4}
\end{equation}
\begin{equation}
  S_\mu
  =
  \frac{\partial\xi^a}{\partial x^\mu}S_a,
  \tag{5.1.5}\label{eq:5.1.5}
\end{equation}
where \(u^a\) and \(S_a\) are components of \(\mathbf u\) and \(\mathbf S\)
in the freely falling coordinate system \(\xi^a\). Although \(u^\mu\) and
\(S_\mu\) are vectors, \(\dd u^\mu/\dd\tau\) and
\(\dd S_\mu/\dd\tau\) are not, but we saw in Section~\ref{sec:4.9} that
there can be defined vector derivatives \(\Dworldline{u^\mu}\) and
\(\Dworldline{S_\mu}\), which reduce when
\(\Gamma^\lambda{}_{\mu\nu}=0\) to the ordinary derivatives
\(\dd u^\mu/\dd\tau\) and \(\dd S_\mu/\dd\tau\). The correct equations for
the particle position and spin are dictated by the Principle of General
Covariance to be
\begin{equation}
  \Dworldline{u^\mu}=0,
  \qquad
  \Dworldline{S_\mu}=0,
  \tag{5.1.6}\label{eq:5.1.6}
\end{equation}
or, in more detail,
\begin{equation}
  \frac{\dd u^\mu}{\dd\tau}
  +
  \Gamma^\mu{}_{\nu\lambda}u^\nu u^\lambda
  =0,
  \tag{5.1.7}\label{eq:5.1.7}
\end{equation}
\begin{equation}
  \frac{\dd S_\mu}{\dd\tau}
  -
  \Gamma^\lambda{}_{\mu\nu}u^\nu S_\lambda
  =0.
  \tag{5.1.8}\label{eq:5.1.8}
\end{equation}
In addition, Eq.~\eqref{eq:5.1.3} becomes now
\begin{equation}
  S_\mu u^\mu=0.
  \tag{5.1.9}\label{eq:5.1.9}
\end{equation}

To repeat the reasoning of Section~\ref{sec:4.1}, these equations are true in
the presence of gravitational fields because they are generally covariant and
true in the absence of gravitation, reducing when
\(\Gamma^\lambda{}_{\mu\nu}\) vanishes to
Eqs.~\eqref{eq:5.1.1}--\eqref{eq:5.1.3}. That is, the Principle of
Equivalence tells us that there are locally inertial coordinate systems in
which Eqs.~\eqref{eq:5.1.6}--\eqref{eq:5.1.9} are valid (provided always
that our particle is sufficiently small), and general covariance then ensures
that these equations hold in the laboratory reference frame.

We recognize in Eqs.~\eqref{eq:5.1.7} and \eqref{eq:5.1.8} the differential
equations for parallel transport of the vectors \(u^\mu\) and \(S_\mu\).
Since \(u^\mu=\dd x^\mu/\dd\tau\), Eq.~\eqref{eq:5.1.7} is nothing but the
familiar equation for free fall, derived previously by differentiating
Eq.~\eqref{eq:5.1.4} with respect to \(\tau\) and using
Eq.~\eqref{eq:5.1.1}; it should be evident that a good deal of work is saved
by using general covariance instead of our previous direct approach.
Equation~\eqref{eq:5.1.8} describes the precession of gyroscopes in free
fall, and is further discussed in Chapter~9. For the present, we note that
\(S_\mu S^\mu\) is constant, because the ordinary derivative of a scalar is
the same as its covariant derivative:
\begin{equation}
  \frac{\dd}{\dd\tau}\left(S_\mu S^\mu\right)
  =
  \frac{D}{D\tau}\left(S_\mu S^\mu\right)
  =0.
  \tag{5.1.10}\label{eq:5.1.10}
\end{equation}

If the particle is not in free fall, then \(\Dworldline{u^\mu}\) does not
vanish, and instead of Eq.~\eqref{eq:5.1.7} we have
\begin{equation}
  \Dworldline{u^\mu}
  \equiv
  \frac{f^\mu}{m},
  \tag{5.1.11}\label{eq:5.1.11}
\end{equation}
where \(m\) is the particle mass and \(f^\mu\) is a contravariant force
vector. This can also be written as
\[
  m\frac{\dd^2x^\mu}{\dd\tau^2}
  =
  f^\mu
  -
  m\Gamma^\mu{}_{\nu\lambda}
  \frac{\dd x^\nu}{\dd\tau}
  \frac{\dd x^\lambda}{\dd\tau}.
\]
The term containing \(m\Gamma^\mu{}_{\nu\lambda}\) evidently plays the role
of a gravitational force. We can always calculate \(f^\mu\) if we know its
value \(f_{\mathrm f}^a\) in freely falling frames \(\xi^a\), for the
requirement that \(f^\mu\) behave as a vector gives it uniquely as
\begin{equation}
  f^\mu
  =
  \frac{\partial x^\mu}{\partial\xi^a}f_{\mathrm f}^a.
  \tag{5.1.12}\label{eq:5.1.12}
\end{equation}
The electromagnetic force will be determined in the next section.

It sometimes happens that a particle is acted on by a force \(f^\mu\)
without experiencing any torque. In this event, an observer in a locally
inertial coordinate frame that is momentarily at rest with respect to the
particle will see no precession of the spin axis; that is,
\(\dd\mathbf S/\dd t\) will vanish. But in this particular coordinate system
\(\dd\mathbf x/\dd t\) also vanishes, so we may write the condition of zero
torque as a Lorentz-invariant statement
\[
  \frac{\dd S^a}{\dd\tau}\propto u^a,
\]
and this will be valid in any locally inertial coordinate system, comoving or
not. Now, what is the constant of proportionality? Let us set
\[
  \frac{\dd S^a}{\dd\tau}=\Phi u^a.
\]
We recall that \(S_a\) is defined so that
\[
  S_a u^a=0,
\]
and therefore
\[
  0
  =
  \frac{\dd}{\dd\tau}\left(S_a u^a\right)
  =
  \Phi u_a u^a
  +
  S_a\frac{\dd u^a}{\dd\tau},
\]
so
\[
  \Phi
  =
  S_a\frac{\dd u^a}{\dd\tau}
  =
  S_a\frac{f^a}{m}.
\]
The spin vector therefore suffers a change given by
\begin{equation}
  \frac{\dd S^a}{\dd\tau}
  =
  \left(S_b\frac{f^b}{m}\right)u^a.
  \tag{5.1.13}\label{eq:5.1.13}
\end{equation}
This phenomenon is known as the \textit{Thomas
precession}.\textsuperscript{\hyperref[ch5-ref-1]{1}} If we now turn on a
gravitational field, Eq.~\eqref{eq:5.1.13} and the Principle of General
Covariance tell us that the spin precesses according to the rule
\begin{equation}
  \Dworldline{S^\mu}
  =
  \left(S_\nu\frac{f^\nu}{m}\right)u^\mu
  =
  S_\nu\Dworldline{u^\nu}\,u^\mu.
  \tag{5.1.14}\label{eq:5.1.14}
\end{equation}
A vector obeying this differential equation is said to be defined by
\textit{Fermi transport};\textsuperscript{\hyperref[ch5-ref-2]{2}} parallel
transport is the special case for \(f^\mu=0\).
